\chapter{LITERATURE REVIEW}

\ChapterSection{Digital Agriculture and Crop Disease Diagnosis}

Agriculture remains a core national sector in Nepal and a direct livelihood concern for rural
households. The Government of Nepal's agricultural statistics report records agriculture as a
major component of national production, while the National Sample Census of Agriculture 2021/22
was designed to provide benchmark data on holdings, land use, crops, livestock, irrigation,
machinery, labour, and agricultural inputs across all districts and local levels
\parencite{moaldAgricultureStats2080,nsoAgricultureCensus2021}. This national context matters
for Krishi Vaidya because crop disease diagnosis is not an isolated technical problem. It is a
field-level decision-support problem that affects farmers who must decide whether a visible crop
symptom requires treatment, monitoring, or expert consultation.

Plant disease and pest pressure is also a global food-security issue. FAO reports that up to
40 percent of global crop production can be lost annually due to plant pests and diseases
\parencite{faoPlantHealthLosses}. The exact loss varies by crop, location, season, management
practice, and pathogen pressure, but the underlying implication is consistent: delayed or
incorrect diagnosis can convert a manageable crop-health problem into a production and income
loss. For smallholder systems, this risk is amplified because advisory services, laboratory
diagnosis, and trained plant-health experts are not always available at the time and place where
the farmer first observes the symptom.

\ChapterSubsection{Farmer-facing diagnosis systems}

Farmer-facing diagnosis systems must solve a different problem from laboratory plant pathology.
They must accept imperfect field images, guide a non-specialist user, and return information in a
form that is understandable and actionable. A useful system should therefore combine visual
analysis with farmer-oriented communication. Image recognition alone may identify a visual class,
but an advisory system must also explain the likely issue, severity, possible management steps,
and the uncertainty of the result.

Krishi Vaidya follows this practical direction. The project does not treat model prediction as the
only output. Instead, the mobile app, backend, YOLO detector, VLM diagnosis pipeline, and advisory
response work together so that image capture, analysis, record keeping, and farmer-facing output
are connected within one system. This is consistent with the direction of digital agriculture in
Nepal. FAO's Nepal digital agriculture work highlights the need for stronger digital
infrastructure, skilled users, coordinated strategy, and advisory services that can support rural
livelihoods \parencite{faoNepalDigitalAgriculture}.

\ChapterSubsection{Dataset and field-variability challenges}

Dataset quality is a central challenge in crop disease diagnosis. Public plant-disease datasets
make model development possible, but they do not automatically represent every field condition.
The PlantVillage study demonstrated the value of image-based plant disease recognition using a
large curated leaf-image dataset \parencite{mohanty2016plantvillage}. However, controlled or
semi-controlled leaf imagery can differ substantially from images captured by farmers in real
fields. Field images may contain mixed lighting, partial occlusion, cluttered background, motion
blur, multiple plant organs, and symptoms at different stages.

This limitation explains why Krishi Vaidya uses separate but connected model components. The YOLO
component focuses on locating plant structures relevant to the image, while the VLM component
produces structured disease-oriented interpretation. The project therefore separates object
localization from disease reasoning rather than forcing a single classifier to perform every task.
The literature supports this separation because plant disease analysis can involve classification,
detection, segmentation, and diagnosis, depending on the visual target and the desired output
\parencite{liu2021plantDiseaseReview}.

\ChapterSection{Object Detection in Agricultural Vision}

Object detection is appropriate when the system must identify not only what is present in an image
but also where it appears. In agricultural images, this is useful because disease symptoms and
plant structures may occupy only part of the image. A detector can localize leaves, panicles,
fruits, ears, heads, and other crop parts that may then support downstream reasoning or user
feedback. For Krishi Vaidya, this is the role of the YOLO model: it provides a fast
object-detection layer for plant-part localization rather than replacing the disease-diagnosis
logic.

\ChapterSubsection{YOLO-family detectors}

YOLO-family detectors are widely used where real-time or near-real-time object detection is
important. The Ultralytics YOLOv8 documentation describes YOLOv8 as a model family supporting
tasks such as detection, segmentation, classification, pose estimation, and oriented bounding
boxes, with multiple model sizes and deployment modes \parencite{ultralyticsYolov8}. The relevance
to Krishi Vaidya is practical: a mobile agricultural system benefits from a detector that can be
trained on custom crop classes, evaluated with standard detection metrics, and exported into
deployment-oriented formats.

The YOLO section of this manuscript therefore uses YOLO terminology carefully. The model is not
claimed to diagnose disease. Its manuscript role is object detection: training on the Krishi
Vaidya bouncer dataset, validating localization performance, comparing thresholds, and evaluating
deployment variants. This distinction is important because object detection and disease diagnosis
measure different capabilities.

\ChapterSubsection{Detection metrics}

Object detection performance is normally evaluated through metrics that combine localization and
classification behavior. Precision measures how many predicted detections are correct. Recall
measures how many ground-truth objects are found. F1 score summarizes the balance between
precision and recall. Mean average precision at an intersection-over-union threshold evaluates
detection quality across classes and confidence thresholds. These metrics are useful because an
agricultural detector must avoid two opposite failures: false detections that confuse the
downstream workflow and missed detections that hide relevant plant structures.

Krishi Vaidya uses these metrics in the YOLO validation and results sections. The project also
uses threshold analysis because the default threshold may not be the best deployment threshold for
the intended user experience. A higher threshold may reduce false positives but miss subtle or
small objects; a lower threshold may increase recall but introduce noisy detections. The
manuscript therefore discusses threshold behavior as an implementation decision, not only as a
reported number.

\ChapterSubsection{Deployment and quantization}

Deployment constraints are important in mobile and field-oriented agricultural systems. A model
that performs well on a workstation may still be impractical if it is too large, too slow, or
incompatible with the target runtime. Ultralytics documents model export paths such as ONNX,
TensorFlow Lite, and quantized formats, including INT8-oriented deployment options
\parencite{ultralyticsExport}. This is directly relevant to Krishi Vaidya because the YOLO work
includes deployment evaluation and because the mobile app contains TFLite model assets.

The literature and tooling therefore justify evaluating more than raw accuracy. The correct
question is whether the trained detector can support an end-user workflow under realistic
resource constraints. This manuscript treats quantization as a deployment tradeoff: compression
and runtime feasibility must be evaluated together with the effect on prediction quality.

\ChapterSection{Vision-Language Models for Diagnosis}

Vision-language models extend the disease-diagnosis workflow beyond fixed-label image
classification. A VLM can consume an image and generate structured text, making it suitable for
tasks where the output must include disease name, severity, explanation, and advisory content. In
Krishi Vaidya, the VLM is used for disease-oriented reasoning and structured output, while the
backend and mobile app transform that output into a farmer-facing response.

\ChapterSubsection{Multimodal disease reasoning}

Qwen2-VL and Qwen2.5-VL provide relevant foundations for Krishi Vaidya because they are
vision-language model families designed to process images and text jointly. Qwen2-VL introduced
dynamic-resolution visual processing and multimodal positional encoding to improve perception
across different visual inputs \parencite{wang2024qwen2vl}. Qwen2.5-VL further reports advances
in visual recognition, object localization, document parsing, structured extraction, and
multimodal reasoning \parencite{bai2025qwen25vl}. These capabilities align with the Krishi
Vaidya requirement for image-based disease interpretation and structured outputs.

The project does not treat the VLM as an unrestricted chatbot. The VLM pipeline is constrained by
prompt templates, class taxonomy, JSON-oriented records, prediction parsing, and validation
metrics. This design is necessary because agricultural diagnosis must be reproducible enough for
manuscript evaluation and safe enough for advisory presentation.

\ChapterSubsection{LoRA fine-tuning}

Fine-tuning a full multimodal model can be computationally expensive. LoRA addresses this by
freezing the base model weights and training low-rank adapter matrices, reducing the number of
trainable parameters required for adaptation \parencite{hu2021lora}. This method is suitable for
student and research environments where the goal is task adaptation without the cost of full
model retraining.

Krishi Vaidya uses LoRA-based fine-tuning for the VLM component. The manuscript therefore reports
the base model, adapter configuration, dataset format, and training behavior as method-level
details. This is important because the final diagnosis behavior depends not only on the base VLM
but also on how the adapter was trained and how predictions were parsed.

\ChapterSubsection{Quantized fine-tuning}

Quantized fine-tuning reduces memory requirements further by combining low-bit model loading with
adapter training. QLoRA demonstrates that gradients can be back-propagated through a frozen
4-bit quantized model into LoRA adapters, enabling efficient fine-tuning of large language models
\parencite{dettmers2023qlora}. Krishi Vaidya's VLM section uses this line of work as the basis
for discussing quantized training configuration and resource behavior.

The manuscript is careful not to overstate this result. Quantization improves feasibility, but it
does not guarantee high diagnostic accuracy. The VLM results in Chapter 4 therefore distinguish
between successful training execution, resource feasibility, dataset validity, and the current
prediction-quality limitations observed in the validation outputs.

\ChapterSubsection{Structured prediction and parseability}

For a farmer-facing advisory system, the model output must be parseable. A natural-language answer
that cannot be reliably converted into application fields is difficult to store, validate, or show
consistently in the mobile app. This is why Krishi Vaidya emphasizes structured output and
prediction parsing. The VLM can generate diagnosis-oriented content, but the backend and mobile
layers require stable fields such as crop, disease, confidence, severity, advice, and uncertainty.

The VLM results show that parseability is not a minor detail. A model can complete training while
still producing a large number of outputs that fail the expected parser. The literature review
therefore supports a central design claim of this manuscript: multimodal diagnosis systems must
evaluate both semantic correctness and output reliability.

\ChapterSection{Mobile and Backend Systems for Agricultural Advisory}

The backend and mobile app convert model outputs into an operational system. Without these
layers, the trained models would remain isolated experiments. Krishi Vaidya uses an API-centered
backend for authentication, scan handling, data persistence, diagnosis orchestration, advisory
generation, and user records. The mobile app provides the farmer-facing interface for image
capture, scan submission, crop management, local storage, synchronization, and diagnosis display.

\ChapterSubsection{Mobile scanning workflows}

React Native is an appropriate implementation basis for cross-platform mobile development because
it enables native Android and iOS applications using React and platform-native capabilities
\parencite{reactNativeDocs}. Expo adds a structured development workflow, routing support, and
tooling for building and deploying React Native applications \parencite{expoDocs}. In Krishi
Vaidya, this technology choice supports rapid development of the scan workflow, crop records,
local database, localization, and backend communication.

The agricultural requirement shapes the mobile design. A scan workflow should not assume perfect
connectivity or expert users. It must guide image selection, preserve records, handle failures,
and present output in a way that is understandable to the farmer. This is why the mobile
implementation includes repository/provider structure, local persistence, sync status, and
farmer-facing result screens rather than acting only as a thin image-upload client.

\ChapterSubsection{Offline-first and synchronization patterns}

Rural and field environments often include unstable connectivity. Krishi Vaidya therefore treats
offline behavior as a functional requirement. The mobile app stores records locally using SQLite
and uses synchronization logic to reconcile local and remote state when connectivity permits. The
literature on mobile systems commonly frames this as an offline-first problem: the app should
remain usable when the network is unavailable, and synchronization should occur later without
losing user data.

The current manuscript does not claim that all offline cases have been proven by device-level
testing. It documents the implemented local storage and synchronization structure, then marks
runtime evidence as a remaining validation requirement. This distinction is necessary for a
high-quality thesis: implementation evidence and field-readiness evidence are not the same.

\ChapterSubsection{API-backed advisory systems}

The backend uses Express.js and Mongoose. Express documentation defines routing as the mechanism
through which application endpoints respond to client requests, and middleware as the request
processing stack that can modify requests, enforce checks, or end the response cycle
\parencite{expressRouting,expressMiddleware}. Mongoose schemas define the shape of MongoDB
documents and support validation, indexes, methods, middleware, and model construction
\parencite{mongooseSchemas}. These ideas map directly to the backend architecture in Krishi
Vaidya: routes expose the API surface, middleware handles security and validation, services
implement business logic, repositories isolate persistence, and Mongoose models define stored
records.

The VLM inference layer uses FastAPI. FastAPI is a Python framework for building APIs using
standard Python type hints \parencite{fastapiDocs}. This fits the VLM deployment boundary because
the model pipeline is Python-based, while the main backend remains a Node.js/Express service. The
result is a two-service architecture: the backend handles application state and orchestration,
while the VLM service handles multimodal inference.

\ChapterSection{Research Gap}

The reviewed literature and project context show three gaps that Krishi Vaidya addresses at a
student-project scale.

First, many plant-disease systems stop at image classification. Such systems can be useful, but
they do not fully address the workflow needed by a farmer who requires image capture, diagnosis,
advice, history, and record keeping. Krishi Vaidya extends beyond isolated classification by
combining object detection, VLM-based interpretation, backend orchestration, and a mobile app.

Second, model performance alone is not enough for deployment-oriented agricultural software. A
usable system must consider model export, structured output, parseability, API contracts,
offline behavior, security, and validation. Krishi Vaidya documents these engineering boundaries
explicitly, including the limitations that remain unresolved.

Third, Nepal's digital agriculture context requires practical, accessible, and mobile-oriented
tools. Official sources identify agriculture as a nationally important sector and digital
agriculture as an area requiring infrastructure, skills, coordination, and advisory-service
development \parencite{moaldAgricultureStats2080,nsoAgricultureCensus2021,faoNepalDigitalAgriculture}.
Krishi Vaidya is positioned within this gap as an integrated prototype for crop-health support,
not as a replacement for expert agronomists or laboratory diagnosis. Its contribution is the
integration of trained model components with a backend and mobile workflow suitable for further
validation, deployment hardening, and field evaluation.
